\section{A separate compiler improvement: the modulus supplies the index}
\label{sec:homogeneous-closure}

The results above improve complete equations in the existing native-compiler
family.  This section concerns a different source architecture.  It removes
one existential coordinate from the inert-prime CRT certificate of source~75.
Its factorial, binomial and variable-exponent operations have not been
compiled and counted as a new ordinary polynomial.  Thus it is a proved
compiler lemma, not another claimed Con(ZFC) endpoint.

\subsection{The homogeneous substitution}

Fix an explicit polynomial
\(R(i,w_1,\ldots,w_m;\mathbf x)\in\mathbb Z[i,\mathbf w,\mathbf x]\),
where \(m\ge1\), and let \(R_{\rm abs}\) replace each coefficient by its
absolute value.  The external parameters \(\mathbf x\) are natural numbers.
For \(n\ge1\) and \(U\ge0\), put
\begin{align*}
 H&=R_{\rm abs}(n,U,\ldots,U;\mathbf x)+n+U+1,&
 s&=4H!,\\
 m_i&=s(i+1)-1\quad(0\le i<n),&
 B&=\prod_{i=0}^{n-1}m_i.
\end{align*}
If an arithmetic DAG already represents \(R=e-f\) with nonnegative
coefficient polynomials, the majorant \(e+f\) can be used instead, preserving
that sharing.  Write \(R=\sum_{a=0}^d R_a(\mathbf w;\mathbf x)i^a\), with
\(d=0\) for the zero polynomial, and define the ordinary integer polynomial
\begin{equation}
 \widehat R(s,\mathbf w;\mathbf x)
 =\sum_{a=0}^{d}R_a(\mathbf w;\mathbf x)(1-s)^a s^{d-a}.
 \label{eq:homogeneous-index}
\end{equation}
The right side is a finite polynomial definition.  The notation
\(s^dR((1-s)/s,\mathbf w;\mathbf x)\) explains the construction, but does
not introduce division into its polynomial representation.

For a fixed binary tree, define a norm recursively by
\(\mathcal N(z)=z\) and
\(\mathcal N(L,R)=\mathcal N(L)^2+\mathcal N(R)^2\).
Use the convention \(\binom Ak=0\) when \(0\le A<k\).

\paragraph{Homogeneous-index closure theorem.}
For each fixed \(n,U,\mathbf x\) as above, the bounded assertion
\[
 \forall i<n\;\exists\mathbf w\in\{0,\ldots,U\}^{m}
                    \quad R(i,\mathbf w;\mathbf x)=0
\]
is equivalent to the existence of natural \(A_1,\ldots,A_m\) such that
\begin{equation}
 B\mid\mathcal N\left(
   \binom{A_1}{U+1},\ldots,\binom{A_m}{U+1},
   \widehat R(s,\mathbf A;\mathbf x)\right).
 \label{eq:homogeneous-certificate}
\end{equation}
There is no separately encoded index and no index-congruence leaf.

\begin{proof}
Every prime divisor of \(m_i\) exceeds \(H\), because a prime at most
\(H\) divides \(H!\) and makes \(m_i\equiv-1\).  The \(m_i\) are pairwise
coprime: a common prime divisor of \(m_i,m_j\) would divide \(j-i\), since
it cannot divide \(s\), but \(0<|j-i|<n\le H\).  Each \(m_i\) is
\(3\pmod4\), so it has a prime divisor \(p_i\equiv3\pmod4\).
Over \(\mathbb F_{p_i}\), a sum of two squares vanishes only when both
entries vanish.  Induction on the norm tree therefore makes vanishing of
its value equivalent to vanishing of all its leaves.

Suppose \eqref{eq:homogeneous-certificate} holds.  For each \(i\), all leaves
vanish modulo \(p_i\).  The identity
\[
 (U+1)!\binom{A_j}{U+1}=\prod_{r=0}^{U}(A_j-r)
\]
and \(p_i>H>U+1\) give \(A_j\equiv w_j\pmod{p_i}\) for some
\(0\le w_j\le U\).  The same prime works for all coordinates.
Furthermore \(s(i+1)\equiv1\pmod{p_i}\), so
\[
 \widehat R(s,\mathbf A;\mathbf x)
       \equiv s^dR(i,\mathbf w;\mathbf x)\pmod{p_i}.
\]
Since \(s\) is invertible modulo \(p_i\), the residual is zero modulo
\(p_i\).  Its absolute value is at most
\(R_{\rm abs}(n,U,\ldots,U;\mathbf x)<H<p_i\), so it is zero in
\(\mathbb Z\).

Conversely, select bounded witnesses for every \(i<n\).  CRT gives
natural \(A_j\) with \(A_j\equiv w_{ij}\pmod{m_i}\).  Every \(m_i\) is
coprime to \((U+1)!\), so all binomial leaves vanish modulo \(m_i\).
The same homogeneous identity and the local equality show that the
residual leaf does also.  Pairwise coprimality then proves
\eqref{eq:homogeneous-certificate}.
\end{proof}

The proof is finite arithmetic formalizable in ZFC; it uses no consistency
or arithmetic-soundness assumption.  Existentially quantifying \(U\)
recovers unbounded witnesses at the finitely many indices.  Since the norm
has at least two leaves, it is nonnegative, and its divisibility can be
written as one equality \(\mathcal N=Bv\) with \(v\in\mathbb N\).

\subsection{The finite product and the exact saving}

As in source~75, the product defining \(B\) has the special-function
expression
\begin{align*}
 Q&=(sn)^n+1,& M&=sQ+1,\\
 B&=s^n n!\binom{Q+n}{n}\bmod M.
\end{align*}
Indeed the numerator is \(\prod_{j=1}^{n}s(Q+j)\), whose factors reduce
to \(sj-1\).  Their positive product is below \((sn)^n<Q<M\), so its
least remainder is the required product.  The variable exponent and the
special functions remain genuine compilation obligations.

Source~75 used the leaves
\[
 \binom{A_j}{U+1}\ (1\le j\le m),\qquad
 R(c,\mathbf A;\mathbf x),\qquad s(c+1)-1.
\]
Equation~\eqref{eq:homogeneous-index} removes the existential coordinate
\(c\), rather than merely dropping its check.  The number of leaves falls
from \(m+2\) to \(m+1\).  A balanced norm tree with \(k\) leaves has maximum
degree \(2^{\lceil\log_2k\rceil}\) in its leaf symbols.  When
\(m=2^r-1\), \(r\ge1\), this bound falls from \(2^{r+1}\) to \(2^r\).
This is not an unconditional degree reduction after substitution and
compilation: homogenization changes the degree in \(s\), and all graph
witnesses and shared operations still need to be counted.

More generally, a modulus of the form \(si+t\), with \(s\) a unit modulo
the selected prime divisor, supplies the residue \(i=-t/s\).  The
corresponding finite homogeneous polynomial uses \((-t)^a s^{d-a}\).
This is a reusable way of recovering an index from the modulus itself.

\subsection{Checks and limits}

The authored standard-library checker verifies 1,125 exact homogeneous
identities, 650 modular index identities, 35 product and coprimality cases,
156,584 finite-field norm evaluations, and 896 full CRT fixtures.  Of the
last set, 79 satisfy the certificate, 594 fail a range condition, and 223
fail the local residual.  Four negative controls check why one needs an
inert prime, the nested norm, a residual bound and an invertible scale.
The tests run with ordinary and optimized Python and contain no assertions
that optimization could suppress.  The universal theorem rests on the
written proof, not these finite diagnostics.

This component does not instantiate a new proof-checker source, eliminate
the special-function graphs, or establish a new minimum-variable
consistency equation independent of K215.  Its exact contribution is the
one-coordinate, one-leaf improvement to a separately identified compiler
stage.
